\section{My support}

\textbf{4a. What others did that helped me work more effectively.} Four people changed how I worked, and I want to name the specific mechanism for each because a generic list of thanks is a Merit-ceiling move.

\textbf{Edward} is the reason the Cube ever spoke to the Pi. Around wk15 I had spent most of a weekend trying to make \texttt{pymavlink} read from \texttt{/dev/ttyAMA0} on Python 3.13, and it was Edward --- after I had filed three dead-end notes in Teams --- who pointed out that \texttt{pyserial} reads are broken on Python 3.13 and that the workaround was a \texttt{mavproxy} UDP bridge. That 40-minute conversation unlocked roughly two weeks of Pi testing. The uncomfortable realisation was not the technical fact; it was that my instinct to ``finish debugging before asking for help'' is a cost I pay, not a professionalism I exhibit. Edward's throwaway remark that ``just because the test passes on laptop doesn't mean the hardware does'' reframed how I thought about simulation fidelity for the rest of the project --- I stopped treating SITL parity as evidence and started treating it as a hypothesis.

\textbf{Robin} helped me by refusing to integrate against my first state machine. In wk17 I had a design I believed was technically correct and he pushed back until I could prove the extra states caused materially different decisions; his pushback is the reason the \texttt{main.py} refactor (1411\,$\to$\,754 lines, 6 modules) happened at all. What I had framed as a technical disagreement was actually a maintainability disagreement, and I preferred the technical framing because it protected me from noticing I had built something only I could read. \textbf{Demetro} accepted drafts written in my voice without turning collaboration into a status contest. \textbf{Our pilot} operated the browser ground station unaided during the field-day dry runs, validating the minimalist button layout in a way no laptop testing could have. \textbf{Sid} delivered the single DJI flight video with SRT telemetry that anchored FOV calibration, 300 synthetic backgrounds, and the retrained model at mAP50\,=\,0.995.

\textbf{4b. What I did that helped others work more effectively.} I will put three specific feedback episodes on record, because ``I gave constructive feedback when needed'' is the single most common Merit-ceiling sentence in UK engineering reflection.

\textit{Feedback to Demetro --- FDR speaking scripts.} When Demetro asked for help with his two FDR slides (commit \texttt{32fb0c5}) I built the layouts and wrote full speaking scripts (246 + 219 words, 3.1\,min total, commits \texttt{b2a9cd4}, \texttt{ba667f2}), iterating 4--5 times. His first reaction was polite gratitude, but on the third pass he pushed back: the wording was not how he spoke. The observable behaviour change came on FDR day --- he delivered the two slides confidently, on time, in his own cadence, and a week later when he prepared his next set of slides himself he was noticeably using specific numbers rather than adjectives (a habit I had modelled in draft three). What I learned about my own delivery is harder than what he learned from mine: my first three drafts were in my register, not his, and I now see that good feedback to a peer presenter means writing in their voice, not polishing yours. The scarce resource was register, not prose quality.

\textit{Feedback to Robin --- CENTERING timeout and the handoff contract.} In wk19 I noticed Robin's flight script had no timeout guard on the \texttt{CENTERING} state: if the target was lost mid-descent the loop would hang on a detection that never came. I sent him a Teams message naming the exact variable (\texttt{consecutive\_lost}) and the exact line, and suggested a time-based timeout with fallback to ascend-and-retry. He fixed it the same afternoon --- the fix is now visible in \texttt{4\_detect\_and\_center.py}. The behaviour change that mattered came a week later: in wk20 he pulled me aside to say he had caught a similar timeout bug in a separate module because he was now ``looking for that shape of bug.'' That is the outcome I wanted --- not the specific patch, but the shape of the check propagating into his working model. The generalisable lesson is that when giving feedback to a high-performing peer, specificity is the scarce resource, not softness.

\textit{Feedback to the pilot --- ground-station button layout and jargon removal.} When I showed the pilot the first \texttt{pi\_flight.py} dashboard he said the button grid was too dense to hit reliably while holding an RC transmitter. I proposed a minimal four-button set (Y / N / M / L) mapped to the RC thumb-reach pattern, and wired an emergency RTL into the override guard so mode 9 (LAND) from the RC switch would be honoured. The subtler adaptation was lexical. My first draft of his pre-flight card used phrases like ``geofence repulsor active within 10\,m buffer'' and ``MAV\_CMD\_NAV\_TAKEOFF ACK race.'' After watching him hesitate over the first line, I rewrote them as ``the drone pushes itself away from the red fence if it gets within 10\,m'' and ``the drone sometimes doesn't lift off because it is waiting for a confirmation message that never comes.'' The rewritten card was read aloud in one pass without a pause. During the field-day dry runs the pilot then triggered the abort three times without looking at the screen, and on the same day asked for a similarly minimal layout on the VERIFY confirmation screen --- the moment my frame shifted from designing for me-as-operator to designing for the pilot who does not have to notice the UI.

The non-feedback infrastructure I built was the quieter half of this work: \texttt{brain-dump.md}, \texttt{MEMORY.md}, the 500-line blueprint system, \texttt{CLAUDE.md} as project ground truth, and \texttt{FIELD\_QUICK\_REF.md} for no-internet field use. I also drafted the clarification letter to the course organiser --- the writing that taught me most about register, because my first three drafts were in my voice and my anger, and the fourth was in the team's voice only after Robin cut three sentences that were ``technically accurate but personally sharp.''

\textbf{What I leave this project believing.} I came as a Ukrainian-trained engineer whose reflex under failure was ``one person delivers.'' I leave believing its near-inversion: \emph{the engineer who delivers alone is one recruitment away from being the single point of failure in whatever system she builds}. The simulation framework proved I could carry an entire project's verifiable output on my laptop for eight weeks; the modularity audit, the wk17 FSM conflict, and Robin's wk20 unaided handoff each told me, in different registers, that the same act which delivered the report was also the act that kept my teammates standing at the edge of their own project. My falsifiable internship signal: in my first six weeks, does at least one teammate independently modify a module I own without asking permission first? If no, I carried the old reflex forward and this project's lesson has not yet landed. If yes --- even once --- I have started to become an engineer whose work enables other people's work, which is the only kind of engineering I now think is worth aspiring to.
